\Aufgabe{Personen im Internet}{

Im Internet findet man über jede:n von uns viele Informationen. Manche sind harmlos, andere können in den falschen Händen gefährlich werden.

\begin{itemize}
    \item Recherchiert möglichst viel zu einer Person im Internet und versucht vor allem Informationen zu finden, die der Person unangenehm sein könnten oder die schwer zu finden sind.
    \begin{itemize}
        \item Beispiele: Frau Duschka, Herr Herrmann, Friedrich Merz, Dominik Krause, eure Eltern, ...
    \end{itemize}
    \item Erstellt einen möglich detailreichen Lebenslauf (keine Maximallänge) mit BYCS-Word
    \item Ergänzt anschließend, was man mit böser Absicht mit diesen Informationen negatives machen könnte.\\
    \emph{Zur Inspiration: Identitätsdiebstahl, materieller Diebstahl, Rufschädigung, etc.}
\end{itemize}

\Large \emphColB{Fazit:} \normalsize Haltet anschließend fest, welche Lehren ihr aus dieser Aufgabe zieht.

\LoesungLine{Ergebnisse sind individuell und werden präsentiert.}{4}

}